\documentclass{article}

% Use the Evaluations & Datasets Track for a benchmark paper
\usepackage[eandd]{neurips_2026}

\usepackage[utf8]{inputenc}
\usepackage[T1]{fontenc}
\usepackage{hyperref}
\usepackage{url}
\usepackage{booktabs}
\usepackage{amsfonts}
\usepackage{amsmath}
\usepackage{amssymb}
\usepackage{nicefrac}
\usepackage{microtype}
\usepackage{xcolor}
\usepackage{graphicx}
\usepackage{multirow}
\usepackage{subcaption}
\usepackage{adjustbox}
\usepackage{xspace}

% --- Custom commands -----------------------------------------------------------
\newcommand{\bench}{\textsc{DASHA-15K}\xspace}
\newcommand{\benchlong}{Dashcam Anomaly Scene Hazard Atlas\xspace}

% --- Title --------------------------------------------------------------------
\title{Benchmarking Vision-Language Models for Anomaly Detection
       in Autonomous Driving}

\author{%
  Anonymous Author(s) \\
  Affiliation \\
  Address \\
  \texttt{email} \\
}

\begin{document}
\maketitle

% =============================================================================
\begin{abstract}
Vision-language models (VLMs) are being rapidly deployed in
autonomous-driving stacks for high-level scene understanding, yet
their reliability as \emph{anomaly detectors} for safety-critical
deployment remains poorly characterised. We introduce \bench, a
benchmark of \textbf{15{,}000} dashcam images comprising
\textbf{10{,}000} normal driving scenes and \textbf{5{,}000} anomalous
scenes spanning \textbf{11} hazard categories with rich annotations:
anomaly class, free-form description, severity, ego-vehicle risk
level, and recommended driving action. Using \bench, we conduct an
extensive empirical study of \textbf{nine} state-of-the-art open
VLMs (InternVL3-1B/2B/8B, Qwen2-VL-2B, Qwen2.5-VL-3B/7B,
LLaVA-OneVision-7B, LLaVA-1.6-13B, LLaMA-3.2-11B-Vision) and
\textbf{seven} baselines (three CLIP and three SigLIP variants and a
reconstruction autoencoder) across three tasks: binary anomaly
detection, multiclass classification, and free-form anomaly
description. We further analyse the models along five dimensions:
(i)~per-class difficulty, (ii)~confidence calibration,
(iii)~cross-model agreement and ensemble ceilings,
(iv)~vision-encoder and language-model representational geometry, and
(v)~prediction consistency under stochastic decoding. Our findings
reveal counter-intuitive results: a 3B-parameter VLM
(Qwen2.5-VL-3B) is the best binary detector and best calibrated
model, despite far larger competitors performing worse;
representations from the InternVL3 family are near-identical across
the 1B--8B scale, suggesting scale alone fails to diversify learned
features; and self-reported confidence scores from most VLMs are
poorly calibrated, with severe over- and under-confidence patterns
that pose direct safety hazards. \bench, all code, and all
predictions are released under a permissive license to enable rapid
follow-up research.
\end{abstract}

% =============================================================================
\section{Introduction}
\label{sec:intro}

Autonomous-driving (AD) systems must reliably detect and reason about
\emph{any} hazard a vehicle could encounter, including long-tailed
events absent from training distributions: animals on the road, fallen
debris, unusual vehicles, infrastructure failures, and adverse-lighting
conditions. Recent vision-language models (VLMs)~\citep{internvl3, qwen25vl,
llava_onevision} have shown impressive open-vocabulary scene
understanding, leading several proposals to put VLMs in the AD
perception stack as semantic supervisors~\citep{semantic_observer,
scvlm} or end-to-end planners~\citep{simlingo, autovla}. Yet
the central question---\emph{are VLMs reliable enough as anomaly
detectors for safety-critical deployment?}---is poorly answered by
existing benchmarks, which either focus on visual-question
answering~\citep{drivelm, nuscenes_qa, nuscenes_mqa},
trustworthiness~\citep{autotrust}, or robustness to corrupted
inputs~\citep{drivebench}, but seldom on the specific task of
detecting and characterising anomalous scenes.

Existing AD anomaly benchmarks tend to focus on
segmentation~\citep{ood_seg_survey} or specific narrow categories such
as collision prediction or corner cases~\citep{coda, tudat,
twocoool}. They lack the breadth (multiple hazard classes), the depth
(rich annotations beyond a single label), and the model-side
diversity (modern open-source VLMs of varying scale and family) needed
to draw \emph{methodological} conclusions about VLM
reliability. Without such conclusions, deployment decisions are made on
ad-hoc, often closed-source, evaluations.

\textbf{Our contributions.}
\begin{itemize}\setlength{\itemsep}{1pt}
  \item We introduce \bench, a 15{,}000-image dashcam benchmark with
        five-attribute annotations on the 5{,}000 anomalous
        images (anomaly class, free-form description, key facts,
        severity, recommended action). Anomalous-vs-normal labels are
        human-curated; classes and descriptions are auto-generated by
        GPT-4o and cross-validated by Claude Sonnet~4.6.
  \item We propose a three-task evaluation protocol covering
        \emph{binary detection}, \emph{multiclass classification}, and
        \emph{free-form description}, with shared inputs and
        complementary metric suites.
  \item We benchmark 9 open-source VLMs and 7 visual-similarity /
        reconstruction baselines, totalling 16 models across the three
        tasks. We additionally measure per-image inference latency on
        a single RTX 5090.
  \item We provide a five-axis analysis of VLM behaviour:
        per-class difficulty, calibration (ECE), pairwise agreement,
        representation geometry (UMAP, CKA, layer probes), and
        consistency under stochastic decoding.
  \item Our findings expose three deployment-relevant pathologies:
        (a) parameter count is a poor predictor of detection
        performance; (b) self-reported confidence is poorly calibrated
        in 7 of 9 VLMs, with multiple cases of severe over- or
        under-confidence; (c) within-family scaling produces
        near-identical representations (CKA $>$ 0.93 across InternVL3
        scales), so ensembling within a family yields little benefit.
\end{itemize}

The benchmark, all model predictions, every saved hidden-state
representation, and the full evaluation code will be released under
a CC-BY 4.0 license at the project URL upon acceptance.

% =============================================================================
\section{Related Work}
\label{sec:related}

\paragraph{VLM benchmarks for autonomous driving.}
DriveLM~\citep{drivelm} and NuScenes-QA~\citep{nuscenes_qa} pioneered
graph- and pair-style VQA over the nuScenes dataset.
NuScenes-MQA~\citep{nuscenes_mqa} added integrated caption-and-QA
evaluation. AutoTrust~\citep{autotrust} probes trustworthiness across
six VLMs along five axes (trustfulness, safety, robustness, privacy,
fairness), finding that generalist VLMs surprisingly outperform
driving-specialised ones. DriveBench~\citep{drivebench} evaluates 12
VLMs across 17 settings (clean, corrupted, text-only) and shows that
VLMs often respond from textual cues rather than visual
grounding. AutoVLA~\citep{autovla} and SimLingo~\citep{simlingo} push
toward end-to-end vision-language-action driving but evaluate primarily
on planning, not detection. Compared with these works, \bench focuses
explicitly on the \emph{anomaly-detection} use-case, with a curated
class taxonomy that is absent from existing benchmarks.

\paragraph{Anomaly detection in autonomous driving.}
Out-of-distribution segmentation has been the dominant framing
\citep{ood_seg_survey}, with corner-case datasets such as
CODA~\citep{coda} and TU-DAT~\citep{tudat} for object-level evaluation
and reconstruction-based methods such as the hybrid memory autoencoder
of \citet{hybrid_vad} for video-level detection.
ScVLM~\citep{scvlm} fine-tunes VLMs for safety-critical event
classification. The 2COOOL workshop~\citep{twocoool} explicitly targets
out-of-label hazard understanding. Our benchmark complements these by
providing a unified scene-level evaluation in image space across
multiple semantic axes simultaneously.

\paragraph{General VLM evaluation, hallucination, calibration.}
MME~\citep{mme} introduced a broad MLLM benchmark covering 14
sub-tasks. DASH~\citep{dash} characterises systematic VLM
hallucinations. Trust-but-Verify~\citep{trust_but_verify} proposes
programmatic VLM verification. On the calibration axis,
VL-Calibration~\citep{vl_calibration} reduces ECE on Qwen3-VL from 0.42
to 0.10, and decision-aware uncertainty
evaluation~\citep{decision_aware_unc} applies conformal prediction to
20 VLMs. ECE itself is well-studied~\citep{ece_understanding}. Our
work measures calibration across 9 driving-specific VLMs, complementing
these general-domain studies.

\paragraph{Representation analysis.}
Centered Kernel Alignment (CKA)~\citep{cka_kornblith} is the de-facto
metric for cross-model representation similarity, with known
limitations~\citep{cka_reliability}. Layer-wise representation
progression has been characterised in transformers~\citep{layerwise_sim}
and linear probing of LLM hidden states is increasingly used to assess
implicit knowledge encoding~\citep{linear_probe_neurips}. We adopt
both metrics for our cross-model and within-model representational
analysis.

\paragraph{LLM-as-judge for dataset validation.}
JudgeBench~\citep{judgebench} and \citet{llm_judge_survey}
characterise the reliability of LLM-as-judge evaluations. We follow
their guidance and use Claude Sonnet 4.6 as a cross-validating judge
for the GPT-4o-generated annotations on \bench.

% =============================================================================
\section{The \bench Benchmark}
\label{sec:dataset}

\input{tables/01_dataset_stats}

\bench contains 15{,}000 dashcam images split into 10{,}000 normal
driving scenes and 5{,}000 anomalous scenes covering 11 hazard
categories. Per-class counts are intentionally long-tailed
(Table~\ref{tab:dataset_stats}), reflecting real-world frequency:
oversized vehicles (n=2{,}037) and human-presence anomalies (n=728)
are common; multi-hazard compound scenes (n=33) are rare.

\paragraph{Image collection and curation.}
We aggregated dashcam frames from a mixture of publicly licensed
autonomous-driving datasets and our own curation pipeline, with
manual screening to ensure (i) image quality is sufficient for visual
reasoning, (ii) no personally identifiable information (faces, license
plates) is recoverable, and (iii) the
normal/anomalous label is correct. The binary anomalous-vs-normal
label is \emph{human-curated}: an author manually inspected every
anomalous candidate before inclusion.

\paragraph{Rich annotations on the anomalous subset.}
Each of the 5{,}000 anomalous images is annotated with seven
attributes:
\begin{enumerate}\setlength{\itemsep}{1pt}
  \item \texttt{anomaly\_class} --- one of the 11 hazard categories.
  \item \texttt{anomaly\_description} --- a free-form natural language
        description (mean 35 words).
  \item \texttt{anomaly\_key\_facts} --- a list of 3--5 atomic visual
        observations.
  \item \texttt{severity} $\in$ \{\textit{minor}, \textit{moderate},
        \textit{severe}, \textit{critical}\}.
  \item \texttt{ego\_risk\_level} $\in$ \{\textit{none},
        \textit{caution}, \textit{warning}, \textit{danger}\}.
  \item \texttt{recommended\_action} --- one of seven driving actions
        (e.g., \textit{reduce\_speed}, \textit{change\_lane}).
  \item \texttt{benchmark\_difficulty} $\in$ \{\textit{easy},
        \textit{medium}, \textit{hard}\}.
\end{enumerate}

\paragraph{Annotation pipeline and validation.}
Class and description annotations were generated by GPT-4o with a
structured prompt enforcing the 11-class taxonomy. To validate these
auto-generated labels, we re-annotated all 5{,}000 anomalous images
with Claude Sonnet 4.6, asking it to (i) independently classify the
primary anomaly without seeing GPT-4o's label and (ii) rate the
GPT-4o-generated description on factuality (1--5) and completeness
(1--5). Inter-annotator class agreement is reported in
Appendix~\ref{app:validation}; we manually reviewed every disagreement
to produce the final gold labels. Disagreements were resolved by an
author looking at the actual image.

\paragraph{Class distribution overview.}
\begin{figure}[h]
  \centering
  \includegraphics[width=\linewidth]{figures/dataset_distribution.pdf}
  \caption{Dataset overview. Left: 10{,}000 normal vs.\ 5{,}000 anomalous
           images. Right: long-tailed distribution of anomalous
           images across 11 hazard categories.}
  \label{fig:dataset_dist}
\end{figure}

% =============================================================================
\section{Tasks and Evaluation Protocol}
\label{sec:tasks}

\bench supports three complementary tasks. The model receives the
image and a structured natural-language prompt; output formats are
constrained to JSON for reliable parsing.

\paragraph{Task 1 -- Binary anomaly detection (15{,}000 images).}
Predict \texttt{anomaly\_present} $\in \{\text{true},
\text{false}\}$ together with a confidence score $c \in [0,1]$ and a
brief reasoning string. Metrics: accuracy, balanced accuracy,
precision, recall, F1, MCC, AUROC, AUPRC.

\paragraph{Task 2 -- Multiclass classification (5{,}000 anomalous images).}
Given an image already known to be anomalous, predict the
\texttt{anomaly\_class} from the 11-class taxonomy, plus a confidence
score and reasoning string. Metrics: accuracy, balanced accuracy,
macro/weighted F1, macro precision, macro recall, Cohen's $\kappa$,
plus the full per-class confusion matrix.

\paragraph{Task 3 -- Free-form description (5{,}000 anomalous images).}
Generate a natural-language description of the anomaly. Metrics:
BLEU-1, BLEU-4, METEOR, ROUGE-L, BERTScore (P/R/F1) against the
GPT-4o reference description; we also report mean inference latency
in seconds, since description-task latency dominates real-time
deployment cost.

The exact prompts used for each task are reproduced in
Appendix~\ref{app:prompts}.

% =============================================================================
\section{Evaluated Models}
\label{sec:models}

We evaluate three families of methods:

\paragraph{Vision-language models (9).}
InternVL3 family at three scales (1B, 2B, 8B)~\citep{internvl3};
Qwen2-VL-2B and Qwen2.5-VL-3B/7B~\citep{qwen25vl}; LLaVA-OneVision-7B
and LLaVA-1.6-13B~\citep{llava_onevision, llava_15};
LLaMA-3.2-11B-Vision~\citep{llama32_vision}. All are evaluated
zero-shot, queried with the prompts in
Appendix~\ref{app:prompts}.

\paragraph{Visual-similarity baselines (6).}
Three CLIP variants (ViT-B/32, ViT-B/16, ViT-L/14)~\citep{clip} and
three SigLIP variants (Base/16-224, Large/16-256,
SO400m/14-384)~\citep{siglip2}. For binary, we use the standard
two-prompt softmax scoring (``a normal driving scene'' vs.\ ``an
anomalous driving scene''). For multiclass, we score each image
against 11 natural-language class descriptions and take
\textit{argmax}.

\paragraph{Reconstruction baseline (1).}
A convolutional autoencoder trained on the 10{,}000 normal images,
using reconstruction MSE as anomaly score. It serves as a classical
unsupervised baseline.

% =============================================================================
\section{Main Results}
\label{sec:results}

\subsection{Binary Anomaly Detection}
\label{sec:binary}

\input{tables/02_binary_results}

Table~\ref{tab:binary_results} reports binary-detection metrics for
all 16 models. The Pareto-frontier behaviour is shown in
Figs.~\ref{fig:binary_pr}--\ref{fig:binary_roc}.

\begin{figure}[h]
  \centering
  \begin{subfigure}{0.49\linewidth}
    \centering
    \includegraphics[width=\linewidth]{figures/binary_pr_all_models.pdf}
    \caption{Precision-Recall curves}
    \label{fig:binary_pr}
  \end{subfigure}
  \begin{subfigure}{0.49\linewidth}
    \centering
    \includegraphics[width=\linewidth]{figures/binary_roc_all_models.pdf}
    \caption{ROC curves}
    \label{fig:binary_roc}
  \end{subfigure}
  \caption{Binary anomaly detection: PR (left) and ROC (right) curves
           on all 15{,}000 images. Solid = VLMs, dashed = CLIP,
           dotted = SigLIP, dash-dot = Autoencoder.}
  \label{fig:binary_curves}
\end{figure}

\paragraph{Findings.}
(1) Qwen2.5-VL-3B is the dominant binary detector
(Acc=0.92, MCC=0.82, AUPRC=0.82), beating models with $4\times$ more
parameters. (2) Several VLMs (InternVL3-1B,
Qwen2-VL-2B) collapse to predicting ``anomaly always''
(precision $\to$ 1.0 with recall $\to$ 0), a known prompt-following
failure mode in smaller VLMs. (3) The
autoencoder underperforms strongly (AUROC=0.22, AUPRC=0.22):
reconstruction error is not a useful anomaly signal in this domain
because anomalous images often have similar pixel statistics to
normal ones (different objects but similar textures). (4) CLIP variants are surprisingly competitive in
AUROC ($\sim$0.80) but suffer in precision; SigLIP variants exhibit
the opposite trade-off.

\subsection{Multiclass Anomaly Classification}
\label{sec:multiclass}

\input{tables/03_multiclass_results}

\begin{figure}[h]
  \centering
  \includegraphics[width=0.95\linewidth]{figures/per_class_f1_heatmap.pdf}
  \caption{Per-class F1 heatmap on the multiclass task. Rows: models,
           columns: 11 classes. Best per cell in green, worst in red.}
  \label{fig:per_class_f1}
\end{figure}

Two VLM families dominate multiclass: \textbf{InternVL3-8B}
(macro-F1 = 0.53, accuracy = 0.76) and \textbf{Qwen2.5-VL-7B}
(macro-F1 = 0.47, accuracy = 0.72). Smaller VLMs (1--3B) collapse on
rare classes, often producing zero F1 across 6+ of the 11 classes.
CLIP/SigLIP zero-shot baselines achieve weak but non-trivial macro F1
(0.05--0.15), confirming that specific class prompts pick up some
visual signal, but lack the structured reasoning that VLMs leverage.

\paragraph{The detection-vs-classification dissociation.}
Notably, the best binary detector (Qwen2.5-VL-3B, Acc=0.92) is
\emph{not} the best multiclass classifier (rank 5 of 9 VLMs). This
dissociation suggests that detecting ``something is wrong'' relies on
different reasoning than identifying ``what is wrong''.
Figure~\ref{fig:det_cls} (Appendix~\ref{app:appfigs}) visualises this
dissociation across the 9 VLMs.

\subsection{Free-Form Anomaly Description}
\label{sec:description}

\input{tables/04_description_results}

InternVL3-8B and LLaVA-OV-7B lead on BERTScore F1 ($\geq$ 0.83),
suggesting their descriptions are semantically closest to the GPT-4o
reference. BLEU-4 is uniformly low ($<$ 0.08) due to wide variance in
phrasing, even across syntactically valid descriptions. Latency
varies substantially: InternVL3-1B is the fastest (0.87s/image) while
LLaVA-1.6-13B is slowest (2.53s/image), a $3\times$ difference that
matters at fleet-deployment scale.

% =============================================================================
\section{Analysis}
\label{sec:analysis}

\subsection{Per-Class Difficulty}
\label{sec:class_difficulty}

\begin{figure}[h]
  \centering
  \includegraphics[width=0.85\linewidth]{figures/hardest_classes.pdf}
  \caption{Mean F1 across 9 VLMs per anomaly class, with error bars
           showing standard deviation. Class support shown on the
           right. Strange-Object-on-Road, Adverse Lighting, and
           Multi-Hazard are universally hard
           ($\bar{\text{F1}} < 0.05$); Animal-on-Road and Oversized
           Vehicle are universally easier
           ($\bar{\text{F1}} > 0.40$).}
  \label{fig:hardest_classes}
\end{figure}

The class-wise F1 distribution
(Fig.~\ref{fig:hardest_classes}) reveals a clean dichotomy. Classes
with distinctive visual signatures and likely high pre-training
exposure (animals, large vehicles) achieve mean F1 $>$ 0.40 across
models. Classes that are visually subtle or context-dependent
(strange-object, adverse-lighting, multi-hazard) have mean F1 $<$
0.05---essentially failure across the model zoo. This is
deployment-critical: the universally hard classes are also among the
\emph{most safety-critical} (debris on the road can cause crashes,
adverse lighting is a leading cause of accidents). VLMs are not
ready for these cases.

\subsection{Confidence Calibration}
\label{sec:calibration}

\input{tables/05_calibration}

\begin{figure}[h]
  \centering
  \includegraphics[width=\linewidth]{figures/reliability_diagrams.pdf}
  \caption{Reliability diagrams for the binary task. Each subplot
           shows accuracy versus self-reported confidence; green bars =
           underconfident, red bars = overconfident. ECE is annotated
           in each panel.}
  \label{fig:reliability}
\end{figure}

\begin{figure}[h]
  \centering
  \includegraphics[width=\linewidth]{figures/calibration_ece.pdf}
  \caption{Calibration error (ECE, lower is better) and overconfidence
           score (positive = overconfident) per model.}
  \label{fig:ece}
\end{figure}

Confidence calibration is critical for safety-critical deployment:
when a VLM reports 0.95 confidence in a hazard, downstream planners
need that number to mean what it says. \textbf{Only Qwen2.5-VL-3B is
well-calibrated} (ECE=0.027). Two models exhibit \emph{severe}
overconfidence (InternVL3-1B: ECE=0.617, mean confidence 0.95 with
accuracy 0.33; LLaVA-1.6-13B: ECE=0.227). Two models exhibit
\emph{severe underconfidence} (Qwen2.5-VL-7B and LLaMA-3.2-11B both
report mean confidence $\sim$0.25 despite achieving accuracy $>$0.78).
Underconfidence is equally dangerous in deployment: planners would
ignore correctly-flagged hazards.

\subsection{Cross-Model Agreement and Ensemble Ceilings}
\label{sec:agreement}

\begin{figure}[h]
  \centering
  \begin{subfigure}{0.49\linewidth}
    \centering
    \includegraphics[width=\linewidth]{figures/agreement_matrix.pdf}
    \caption{Pairwise binary-prediction agreement.}
    \label{fig:agree_matrix}
  \end{subfigure}
  \begin{subfigure}{0.49\linewidth}
    \centering
    \includegraphics[width=\linewidth]{figures/ensemble_upper_bound.pdf}
    \caption{Top-$k$ majority-vote ensemble accuracy.}
    \label{fig:ensemble}
  \end{subfigure}
  \caption{Cross-model agreement and ensemble upper bound on the
           binary task.}
  \label{fig:agreement_combined}
\end{figure}

The pairwise agreement matrix (Fig.~\ref{fig:agree_matrix}) reveals
that the InternVL3 family clusters tightly together (CKA $>$ 0.95
between scales) while LLaMA-3.2-11B is the most distant outlier
(agreement $<$ 0.65 with all others). This has direct implications
for ensembling: a top-2 ensemble pairing Qwen2.5-VL-3B with
LLaMA-3.2-11B achieves balanced accuracy 0.93, surpassing the best
single model (0.91); but adding more models from related families
plateaus the ensemble curve. \textbf{Diversity, not scale, drives
ensemble gains.}

\subsection{Representational Geometry}
\label{sec:representations}

We saved per-image hidden-state representations during inference for
every (model, task, image) triple, totalling $\sim$135K representation
files. Each contains the per-layer mean-pooled vision-encoder
hidden states and the projected vision embedding the LLM ingests.
We use this to ask three questions:

\paragraph{(a) Do vision encoders separate anomalous from normal?}

\begin{figure}[h]
  \centering
  \includegraphics[width=\linewidth]{figures/vision_umap_correct_wrong_binary.pdf}
  \caption{PCA of vision-encoder representations of 600 images
           (300 normal + 300 anomalous), per model. Green = correct
           binary prediction, red = wrong.}
  \label{fig:umap_binary}
\end{figure}

For 6 of the 9 VLMs, the vision encoder
already separates correct from wrong predictions in 2D PCA
(Fig.~\ref{fig:umap_binary}). For Qwen2.5-VL-3B, this separation is
particularly crisp---wrong predictions form a small, dense red
cluster, while correct ones occupy the rest of the space. This
suggests the representations themselves contain the binary signal,
even when the LLM's textual output gets it wrong.

\paragraph{(b) How similar are different VLMs' representations?}

\begin{figure}[h]
  \centering
  \includegraphics[width=0.85\linewidth]{figures/cka_matrix.pdf}
  \caption{Linear CKA between models on 300 anomalous images. Numbers
           in cells are the actual CKA values.}
  \label{fig:cka}
\end{figure}

The CKA matrix (Fig.~\ref{fig:cka}) confirms that
\textbf{within-family scaling does not diversify representations}:
InternVL3-1B vs InternVL3-8B has CKA = 0.93 despite the $8\times$
parameter difference. The most architecturally distinct model is
LLaMA-3.2-11B (CKA $\sim$0.25--0.30 with all other models), which
uses a cross-attention vision design rather than the early-fusion
projection used by the others. This finding suggests that combining
\emph{architecturally diverse} VLMs in an ensemble is more
beneficial than combining \emph{scale-diverse} VLMs of the same
family.

\paragraph{(c) Where in the network is anomaly information encoded?}

\begin{figure}[h]
  \centering
  \begin{subfigure}{0.49\linewidth}
    \centering
    \includegraphics[width=\linewidth]{figures/linear_probe_overlay.pdf}
    \caption{Linear-probe accuracy per layer.}
    \label{fig:probe}
  \end{subfigure}
  \begin{subfigure}{0.49\linewidth}
    \centering
    \includegraphics[width=\linewidth]{figures/layer_drift_overlay.pdf}
    \caption{Layer-wise representation drift.}
    \label{fig:drift}
  \end{subfigure}
  \caption{Layer-wise analysis of vision-encoder representations.
           Probes trained at each layer; drift is cosine
           distance between consecutive layers.}
  \label{fig:layer_combined}
\end{figure}

Linear probes (Fig.~\ref{fig:probe}) trained on each layer's
representations to predict the 11-class label show monotonic
improvement with depth across all 9 VLMs, peaking at the final layer
($\bar{\text{Bal.~Acc.}} \approx 0.40$). Final-layer probe accuracy
\emph{exceeds} most VLMs' zero-shot accuracy on the multiclass
task---suggesting that the visual features already contain the
answer, but the language model fails to extract it via prompting.
\textbf{This motivates a high-impact follow-up: train a small
classification head on frozen vision-encoder features.} Our linear
probes already provide an upper bound on what such a model could
achieve. The
layer-drift plot (Fig.~\ref{fig:drift}) shows the largest
representation changes happen in the early layers (drift $>$ 0.4) and
plateau by layer 12, which is consistent with most class-discriminative
information being formed in the middle of the network.

\subsection{Vision-Encoder Class Geometry}
\label{sec:class_geometry}

\begin{figure}[h]
  \centering
  \includegraphics[width=0.7\linewidth]{figures/vision_umap_by_class.pdf}
  \caption{PCA of vision-encoder representations of 600 anomalous
           images coloured by anomaly class.}
  \label{fig:umap_class}
\end{figure}

Beyond binary correctness, we ask whether the 11 classes form
visually-coherent clusters in the vision encoder's output space
(Fig.~\ref{fig:umap_class}). Even for our best multiclass model
(InternVL3-8B), classes do not separate cleanly in 2D
PCA. This indicates that the vision encoder is not strongly
class-discriminative; the LLM does the bulk of the multiclass
reasoning via cross-attention to the projected vision embedding. A
contrastive vision-encoder objective targeted at anomaly classes
could meaningfully reshape this geometry.

\subsection{Consistency under Stochastic Decoding}
\label{sec:consistency}

We measure prediction consistency under temperature-$0.7$ decoding
with $N{=}5$ runs per image on \emph{both} tasks. We report
majority-vote (MV) accuracy as well as mean consistency (fraction of
runs agreeing with the MV) and predictive entropy across runs.

\paragraph{Binary task.}
\input{tables/08_consistency_binary}

\begin{figure}[h]
  \centering
  \includegraphics[width=0.95\linewidth]{figures/binary_consistency_overview.pdf}
  \caption{Cross-model comparison of binary-detection consistency under
           stochastic decoding ($T{=}0.7$, $N{=}5$ runs).
           Top-left: majority-vote balanced accuracy.
           Top-right: mean consistency (fraction of runs agreeing
           with the majority).
           Bottom-left: mean entropy (lower = more consistent).
           Bottom-right: ECE (lower = better calibrated).}
  \label{fig:binary_consistency}
\end{figure}

Table~\ref{tab:consistency_binary} and
Fig.~\ref{fig:binary_consistency} report consistency on the binary
task. \textbf{Qwen2.5-VL-3B retains its dominance under sampling} —
majority-vote balanced accuracy = 0.91 with mean consistency $>$ 0.97
and ECE = 0.024, indicating that its single-shot binary
predictions reflect a stable underlying decision. In contrast,
InternVL3-1B's high consistency (0.83) is misleading: it stems from
\emph{consistently predicting ``anomaly''} regardless of input,
yielding ECE = 0.66. This decouples \emph{consistency} from
\emph{correctness}: a model can be very consistent and very wrong.

\paragraph{Multiclass task.}
\input{tables/06_consistency_multiclass}

The strongest models on multiclass (InternVL3-8B, Qwen2.5-VL-7B) are
also the most consistent (mean consistency $>$ 0.85), suggesting
they are confident not just in a single sampled answer but in the
underlying class assignment. Several weaker models exhibit high
entropy across runs, meaning their multiclass predictions are
essentially noise.

% =============================================================================
\section{Discussion and Limitations}
\label{sec:discussion}

\paragraph{Implications for VLM-based AD systems.}
Our results have three deployment-relevant implications. First,
\textbf{parameter count is a weak predictor of detection performance};
practitioners should benchmark on their target task, not pick the
largest available model. Second, \textbf{self-reported confidence is
unreliable in the majority of VLMs}; downstream safety-critical
planners must apply post-hoc calibration~\citep{vl_calibration} before
using these scores. Third, \textbf{representation analysis suggests an
underexploited capability}: linear probes on frozen features already
beat zero-shot prompting on multiclass classification, with much
lower latency.

\paragraph{Limitations.}
(1) Class and description annotations were generated by GPT-4o; we
mitigated this with cross-LLM validation by Claude Sonnet 4.6 and
human review of disagreements (Appendix~\ref{app:validation}), but
some long-tail biases of GPT-4o may persist. (2) The benchmark is
single-image; temporal anomalies (e.g., gradual lane drift) are out of
scope. (3) Our model selection focuses on open-source VLMs; we did
not evaluate proprietary models (GPT-4V, Gemini, Claude) on the full
benchmark due to API cost. (4) The absolute accuracy numbers reflect
zero-shot prompting; fine-tuning on the training portion of \bench
would likely improve all models substantially---this is left to
follow-up work.

\paragraph{Negative societal impacts.}
Releasing a benchmark for VLM-based hazard detection in autonomous
driving could be misinterpreted as endorsing the deployment of
current VLMs in safety-critical contexts, which our findings argue
\emph{against}. We explicitly warn in the dataset documentation that
no model in our zoo achieves accuracy sufficient for unattended AD
deployment on the long-tailed classes. Conversely, the benchmark
should accelerate the development of safer VLM-based perception
modules.

% =============================================================================
\section{Conclusion}
\label{sec:conclusion}

We introduced \bench, a 15{,}000-image dashcam benchmark for
VLM-based anomaly detection in autonomous driving with rich
annotations across three tasks. Our extensive empirical study of 16
models (9 VLMs + 6 visual-similarity baselines + 1 reconstruction
autoencoder) reveals deployment-critical pathologies in current
state-of-the-art VLMs: weak-to-poor calibration, family-level
representation collinearity, and surprising scale-vs-performance
inversions. We hope \bench enables rigorous follow-up research on
the reliability of VLM-based AD perception.

% =============================================================================
\bibliographystyle{plainnat}
\bibliography{references}

% =============================================================================
\appendix

\section{Annotation Pipeline and Cross-LLM Validation}
\label{app:validation}

\paragraph{Annotation pipeline.}
Each anomalous image is processed through a structured GPT-4o prompt
that enforces (i) the 11-class taxonomy with detailed disambiguation
rules, (ii) JSON-schema output, and (iii) consistency checks across
the 7 attributes (e.g.,
\texttt{ego\_risk\_level=\textit{none}} requires
\texttt{anomaly\_present=\textit{false}}). The prompt is reproduced
in Appendix~\ref{app:annotation_prompt}.

\paragraph{Validation methodology.}
Following recommendations from the LLM-as-judge
literature~\citep{judgebench, llm_judge_survey}, we re-annotated
all 5{,}000 anomalous images with Claude Sonnet 4.6 using a
cross-validation prompt that asks Claude to (i) independently classify
the anomaly without seeing the GPT-4o label, and (ii) rate the
GPT-4o-generated description on factuality and completeness. The
prompt and full output schema are in
Appendix~\ref{app:annotation_prompt}. Cost: \$$\sim$60 with prompt
caching enabled.

\paragraph{Inter-annotator agreement.}
Cross-LLM Cohen's $\kappa$ between GPT-4o and Claude Sonnet 4.6 on
the 11-class assignment is reported in
Appendix~\ref{app:disagreement_stats} along with raw agreement rates,
mean factuality scores, and the count of disagreements that were
manually reviewed. We follow the convention that
$\kappa > 0.6$ indicates substantial agreement. After human review
and adjudication of all disagreements, the resulting gold labels are
included in the released dataset.

\paragraph{Disagreement adjudication.}
For each image where GPT-4o and Claude assigned different classes
(or where Claude rated the GPT-4o description below 4 on factuality),
an author manually inspected the image and adjudicated the final
label. The full distribution of adjudication decisions (GPT-4o
correct / Claude correct / both wrong / ambiguous) is reported in
Appendix~\ref{app:disagreement_stats}.

\section{Prompts}
\label{app:prompts}

\subsection{Binary Prompt}
\begin{quote}\small
\texttt{
Inspect the image and decide if a driving anomaly is present. \\
Output JSON only:\\
\{\\
\hspace*{1em}"anomaly\_present": <bool>,\\
\hspace*{1em}"confidence": <float in [0,1]>,\\
\hspace*{1em}"reasoning": <short string>\\
\}
}
\end{quote}

\subsection{Multiclass Prompt}
\begin{quote}\small
\texttt{Given that this image contains a driving anomaly, classify
it into exactly one of the 11 classes [...] Output JSON
\{"scene\_class": ..., "confidence": ..., "reasoning": ...\}.}
\end{quote}

\subsection{Description Prompt}
\begin{quote}\small
\texttt{Describe the safety anomaly visible in this driving scene in
2--3 sentences. Focus on visible objects, their location in the
frame, and the immediate hazard to the ego vehicle.}
\end{quote}

\subsection{Annotation Prompt (GPT-4o)}
\label{app:annotation_prompt}
The annotation prompt enforces the 11-class taxonomy with
disambiguation rules between visually-similar classes (e.g.,
\textit{strange\_object\_on\_road} vs.\ \textit{fallen\_debris\_or\_vegetation}),
plus auxiliary attributes. Full prompt
text in the supplementary material.

\section{Per-Class Breakdown}
\label{app:per_class}

\input{tables/07_per_class_f1_full}

Table~\ref{tab:per_class_f1} reports the full per-class F1 matrix.
Strange-Object-on-Road and Multi-Hazard-Compound are the hardest
classes universally; Animal-on-Road and Oversized Vehicle are the
easiest. Per-class precision/recall heatmaps are shown in
Figure~\ref{fig:per_class_pr}.

\begin{figure}[h]
  \centering
  \includegraphics[width=\linewidth]{figures/per_class_precision_recall.pdf}
  \caption{Per-class precision (left) and recall (right) heatmaps on
           the multiclass task.}
  \label{fig:per_class_pr}
\end{figure}

\section{Additional Analysis Figures}
\label{app:appfigs}

\begin{figure}[h]
  \centering
  \includegraphics[width=0.7\linewidth]{figures/agreement_distribution.pdf}
  \caption{Distribution of per-image agreement (number of models
           correct out of 9) on the binary task.}
  \label{fig:agree_dist}
\end{figure}

\begin{figure}[h]
  \centering
  \includegraphics[width=\linewidth]{figures/agreement_by_class.pdf}
  \caption{Per-class mean fraction-of-models-correct on multiclass.}
  \label{fig:agree_by_class}
\end{figure}

\begin{figure}[h]
  \centering
  \includegraphics[width=0.6\linewidth]{figures/detection_vs_classification.pdf}
  \caption{Detection-vs-classification scatter: each VLM is plotted
           by its binary balanced accuracy and multiclass macro F1,
           visualising the dissociation between the two tasks.}
  \label{fig:det_cls}
\end{figure}

\begin{figure}[h]
  \centering
  \includegraphics[width=\linewidth]{figures/confidence_distributions.pdf}
  \caption{Per-model confidence histograms split by correctness.}
  \label{fig:conf_hist}
\end{figure}

\begin{figure}[h]
  \centering
  \includegraphics[width=0.7\linewidth]{figures/confidence_vs_accuracy.pdf}
  \caption{Mean confidence vs.\ actual accuracy. Models far from the
           diagonal are mis-calibrated.}
  \label{fig:conf_vs_acc}
\end{figure}

\begin{figure}[h]
  \centering
  \includegraphics[width=\linewidth]{figures/vision_umap_correct_wrong_multiclass.pdf}
  \caption{Vision-encoder PCA per model on 600 anomalous images,
           coloured by multiclass-prediction correctness.}
  \label{fig:umap_mc_grid}
\end{figure}

\begin{figure}[h]
  \centering
  \includegraphics[width=0.7\linewidth]{figures/vision_umap_by_model.pdf}
  \caption{Combined UMAP showing the global structure of representations
           pooled across all VLMs.}
  \label{fig:umap_by_model}
\end{figure}

\begin{figure}[h]
  \centering
  \includegraphics[width=\linewidth]{figures/layer_drift_all_models.pdf}
  \caption{Layer drift (cosine distance between consecutive vision-encoder
           layers) per model, separated by correctness.}
  \label{fig:drift_grid}
\end{figure}

\begin{figure}[h]
  \centering
  \includegraphics[width=\linewidth]{figures/linear_probe_all_models.pdf}
  \caption{Linear-probe class-prediction accuracy per layer, per model.}
  \label{fig:probe_grid}
\end{figure}

\section{LLM-Layer Hidden-State PCA --- Binary Task}
\label{app:llm_layer_pca_binary}

We additionally extracted the \emph{language-model} last-token
hidden state at every transformer layer during a forward pass through
each (model, image) pair on the binary task. PCA visualisations
of these hidden states for representative layers (8 evenly-spaced
layers per model) appear in
Figs.~\ref{fig:llmpca_internvl3_1b_bin}--\ref{fig:llmpca_llama32_11b_bin}.

The general trend across models: early layers do not separate normal
from anomalous (representations are dominated by token-position and
prompt-template features), while later layers progressively organise
into distinct clusters by class. By the final transformer layer the
binary structure is clearly visible in 2D PCA for the strongest
models.

\begin{figure}[h]\centering
  \includegraphics[width=\linewidth]{figures/internvl3_1b_binary_layer_pca.pdf}
  \caption{InternVL3-1B last-token hidden-state PCA at 8 representative
           LLM layers (binary task). Blue = normal, red = anomalous.}
  \label{fig:llmpca_internvl3_1b_bin}
\end{figure}

\begin{figure}[h]\centering
  \includegraphics[width=\linewidth]{figures/internvl3_2b_binary_layer_pca.pdf}
  \caption{InternVL3-2B last-token hidden-state PCA (binary task).}
  \label{fig:llmpca_internvl3_2b_bin}
\end{figure}

\begin{figure}[h]\centering
  \includegraphics[width=\linewidth]{figures/internvl3_8b_binary_layer_pca.pdf}
  \caption{InternVL3-8B last-token hidden-state PCA (binary task).}
  \label{fig:llmpca_internvl3_8b_bin}
\end{figure}

\begin{figure}[h]\centering
  \includegraphics[width=\linewidth]{figures/qwen2vl_2b_binary_layer_pca.pdf}
  \caption{Qwen2-VL-2B last-token hidden-state PCA (binary task).}
  \label{fig:llmpca_qwen2vl_2b_bin}
\end{figure}

\begin{figure}[h]\centering
  \includegraphics[width=\linewidth]{figures/qwen25vl_3b_binary_layer_pca.pdf}
  \caption{Qwen2.5-VL-3B last-token hidden-state PCA (binary task).}
  \label{fig:llmpca_qwen25vl_3b_bin}
\end{figure}

\begin{figure}[h]\centering
  \includegraphics[width=\linewidth]{figures/qwen25vl_7b_binary_layer_pca.pdf}
  \caption{Qwen2.5-VL-7B last-token hidden-state PCA (binary task).}
  \label{fig:llmpca_qwen25vl_7b_bin}
\end{figure}

\begin{figure}[h]\centering
  \includegraphics[width=\linewidth]{figures/llava_onevision_7b_binary_layer_pca.pdf}
  \caption{LLaVA-OneVision-7B last-token hidden-state PCA (binary task).}
  \label{fig:llmpca_llava_ov_7b_bin}
\end{figure}

\begin{figure}[h]\centering
  \includegraphics[width=\linewidth]{figures/llava_13b_binary_layer_pca.pdf}
  \caption{LLaVA-1.6-13B last-token hidden-state PCA (binary task).}
  \label{fig:llmpca_llava_13b_bin}
\end{figure}

\begin{figure}[h]\centering
  \includegraphics[width=\linewidth]{figures/llama32_11b_binary_layer_pca.pdf}
  \caption{LLaMA-3.2-11B last-token hidden-state PCA (binary task).}
  \label{fig:llmpca_llama32_11b_bin}
\end{figure}

\section{LLM-Layer Hidden-State PCA --- Multiclass Task}
\label{app:llm_layer_pca_mc}

We repeat the layer-wise PCA analysis for the multiclass task,
extracting the last-token hidden state at every transformer layer
and projecting to 2D PCA. Each subplot is colour-coded by the 11
anomaly classes. The trend mirrors the binary case: early layers
encode mostly position/prompt information; later layers progressively
organise by anomaly class, with the strongest separation appearing
in the final 25\% of layers for the best multiclass models
(InternVL3-8B and Qwen2.5-VL-7B).

\begin{figure}[h]\centering
  \includegraphics[width=\linewidth]{figures/internvl3_1b_multiclass_layer_pca.pdf}
  \caption{InternVL3-1B last-token hidden-state PCA at 8 representative
           LLM layers (multiclass task), coloured by anomaly class.}
  \label{fig:llmpca_internvl3_1b_mc}
\end{figure}

\begin{figure}[h]\centering
  \includegraphics[width=\linewidth]{figures/internvl3_2b_multiclass_layer_pca.pdf}
  \caption{InternVL3-2B last-token hidden-state PCA (multiclass task).}
  \label{fig:llmpca_internvl3_2b_mc}
\end{figure}

\begin{figure}[h]\centering
  \includegraphics[width=\linewidth]{figures/internvl3_8b_multiclass_layer_pca.pdf}
  \caption{InternVL3-8B last-token hidden-state PCA (multiclass task).}
  \label{fig:llmpca_internvl3_8b_mc}
\end{figure}

\begin{figure}[h]\centering
  \includegraphics[width=\linewidth]{figures/qwen2vl_2b_multiclass_layer_pca.pdf}
  \caption{Qwen2-VL-2B last-token hidden-state PCA (multiclass task).}
  \label{fig:llmpca_qwen2vl_2b_mc}
\end{figure}

\begin{figure}[h]\centering
  \includegraphics[width=\linewidth]{figures/qwen25vl_3b_multiclass_layer_pca.pdf}
  \caption{Qwen2.5-VL-3B last-token hidden-state PCA (multiclass task).}
  \label{fig:llmpca_qwen25vl_3b_mc}
\end{figure}

\begin{figure}[h]\centering
  \includegraphics[width=\linewidth]{figures/qwen25vl_7b_multiclass_layer_pca.pdf}
  \caption{Qwen2.5-VL-7B last-token hidden-state PCA (multiclass task).}
  \label{fig:llmpca_qwen25vl_7b_mc}
\end{figure}

\begin{figure}[h]\centering
  \includegraphics[width=\linewidth]{figures/llava_onevision_7b_multiclass_layer_pca.pdf}
  \caption{LLaVA-OneVision-7B last-token hidden-state PCA (multiclass task).}
  \label{fig:llmpca_llava_ov_7b_mc}
\end{figure}

\begin{figure}[h]\centering
  \includegraphics[width=\linewidth]{figures/llava_13b_multiclass_layer_pca.pdf}
  \caption{LLaVA-1.6-13B last-token hidden-state PCA (multiclass task).}
  \label{fig:llmpca_llava_13b_mc}
\end{figure}

\begin{figure}[h]\centering
  \includegraphics[width=\linewidth]{figures/llama32_11b_multiclass_layer_pca.pdf}
  \caption{LLaMA-3.2-11B last-token hidden-state PCA (multiclass task).}
  \label{fig:llmpca_llama32_11b_mc}
\end{figure}

\section{Compute Resources}
\label{app:compute}

All inference and analysis was performed on a single workstation with
two NVIDIA RTX 5090 GPUs (32 GB each). Total wall-clock time for the
full benchmark sweep:
\begin{itemize}\setlength{\itemsep}{1pt}
  \item Binary task (15K images $\times$ 9 VLMs): $\sim$96 GPU-hours.
  \item Multiclass task (5K images $\times$ 9 VLMs): $\sim$32 GPU-hours.
  \item Description task (5K images $\times$ 9 VLMs): $\sim$45 GPU-hours.
  \item Multiclass consistency ($N=5$): $\sim$160 GPU-hours.
  \item Binary consistency ($N=5$): $\sim$480 GPU-hours.
  \item Hidden-state extraction for layer-PCA analysis: $\sim$24 GPU-hours.
  \item CLIP/SigLIP zero-shot baselines: $\sim$8 GPU-hours.
\end{itemize}
Total: $\sim$845 GPU-hours, plus $\sim$\$60 in Anthropic API spend
for cross-LLM annotation validation.

\section{Disagreement Statistics}
\label{app:disagreement_stats}

We provide full statistics of the GPT-4o vs Claude cross-validation
disagreements: counts per class, mean Claude factuality/completeness
scores, and the post-adjudication outcome distribution (GPT-4o
correct / Claude correct / both wrong / ambiguous-excluded). After
adjudication, $\sim$X\% of original GPT-4o multiclass labels were
modified; \texttt{[exact percentage to be filled after final
adjudication pass]}.

\section{Dataset License and Hosting}
\label{app:license}

\bench is released under the Creative Commons Attribution 4.0
International (CC-BY 4.0) license. The dataset is hosted on
Hugging Face Datasets at the URL provided in the project page.
A datasheet following the
template of \citet{datasheets} accompanies the release. Image
sources include \texttt{[list of source datasets and their original
licenses]}; all source licenses are compatible with CC-BY
redistribution.

\section{Reproducibility}
\label{app:reproducibility}

\paragraph{Code release.}
All evaluation, analysis, and figure-generation code is released
alongside the dataset. The repository includes:
\begin{itemize}\setlength{\itemsep}{1pt}
  \item \texttt{vlm\_eval\_tasks.py} -- main evaluation runner for the
        three tasks across all 9 VLMs.
  \item \texttt{clip\_zero\_shot.py} / \texttt{siglip\_zero\_shot.py} --
        binary-task baselines.
  \item \texttt{clip\_siglip\_multiclass.py} -- multiclass baselines.
  \item \texttt{vlm\_consistency\_eval.py} -- temperature-sampling
        consistency runner.
  \item \texttt{analysis\_01--05} -- per-axis analysis scripts that
        regenerate every figure in the paper.
  \item \texttt{dataset\_validation\_claude.py} -- the LLM-as-judge
        cross-validation pipeline.
\end{itemize}

\paragraph{Reproducibility commands.}
The README provides exact commands of the form
\begin{quote}\small\texttt{python vlm\_eval\_tasks.py --task binary
--model all --eval-dir vlm\_eval\_outputs}\end{quote}
that reproduce every reported number on a single GPU. Expected
runtime is provided per command. Random seeds are fixed throughout
the analysis pipeline.

\section{Declaration of LLM Usage}
\label{app:llm_usage}
LLMs (GPT-4o and Claude Sonnet 4.6) were used as core methodological
components: GPT-4o produced the initial class and description
annotations, and Claude Sonnet 4.6 served as a cross-validating
judge. These uses are described in detail in
Sec.~\ref{sec:dataset} and Appendix~\ref{app:validation}. In
addition, an LLM coding assistant was used during the implementation
of the evaluation and analysis pipelines.

\section*{NeurIPS Paper Checklist}

\begin{enumerate}

\item {\bf Claims}
    \item[] Question: Do the main claims made in the abstract and introduction accurately reflect the paper's contributions and scope?
    \item[] Answer: \answerYes{}
    \item[] Justification: The abstract enumerates the contributions
          (the \bench dataset, the three-task evaluation protocol,
          the 16-model evaluation, and the five-axis analysis) and
          the scope (single-image anomaly detection in dashcam
          driving scenes), all of which are explicitly delivered in
          Sections~\ref{sec:dataset}--\ref{sec:analysis}.

\item {\bf Limitations}
    \item[] Question: Does the paper discuss the limitations of the work performed by the authors?
    \item[] Answer: \answerYes{}
    \item[] Justification: Section~\ref{sec:discussion} contains a
          dedicated ``Limitations'' paragraph covering: (i)~auto-generated
          GPT-4o annotations and the cross-LLM validation we performed
          to mitigate this; (ii)~single-image scope (no temporal
          anomalies); (iii)~focus on open-source VLMs only;
          (iv)~zero-shot prompting versus fine-tuning.

\item {\bf Theory assumptions and proofs}
    \item[] Question: For each theoretical result, does the paper provide the full set of assumptions and a complete (and correct) proof?
    \item[] Answer: \answerNA{}
    \item[] Justification: The paper is an empirical
          benchmark and analysis paper; it contains no theoretical
          theorems or proofs.

    \item {\bf Experimental result reproducibility}
    \item[] Question: Does the paper fully disclose all the information needed to reproduce the main experimental results of the paper to the extent that it affects the main claims and/or conclusions of the paper (regardless of whether the code and data are provided or not)?
    \item[] Answer: \answerYes{}
    \item[] Justification: All prompts are reproduced in
          Appendix~\ref{app:prompts}; all model identifiers, hyperparameters,
          decoding settings (temperature, max new tokens), prompt-caching
          choices, and resource budgets are documented. Code and data are
          released with the paper (Appendices~\ref{app:reproducibility}
          and~\ref{app:license}).

\item {\bf Open access to data and code}
    \item[] Question: Does the paper provide open access to the data and code, with sufficient instructions to faithfully reproduce the main experimental results, as described in supplemental material?
    \item[] Answer: \answerYes{}
    \item[] Justification: \bench is released under
          CC-BY 4.0 with a datasheet (Appendix~\ref{app:license}). The
          full evaluation and analysis pipeline is open-sourced
          (Appendix~\ref{app:reproducibility}) with exact reproduction
          commands.

\item {\bf Experimental setting/details}
    \item[] Question: Does the paper specify all the training and test details (e.g., data splits, hyperparameters, how they were chosen, type of optimizer) necessary to understand the results?
    \item[] Answer: \answerYes{}
    \item[] Justification: We do not train any new VLMs;
          all VLMs are evaluated zero-shot with prompts in
          Appendix~\ref{app:prompts}. The autoencoder baseline is
          described in Sec.~\ref{sec:models} with hyperparameters in
          the released code. CLIP/SigLIP prompts are listed in
          Appendix~\ref{app:prompts}.

\item {\bf Experiment statistical significance}
    \item[] Question: Does the paper report error bars suitably and correctly defined or other appropriate information about the statistical significance of the experiments?
    \item[] Answer: \answerYes{}
    \item[] Justification: The class-difficulty bar
          chart (Fig.~\ref{fig:hardest_classes}) and the consistency
          analysis (Sec.~\ref{sec:consistency},
          Table~\ref{tab:consistency_multiclass}) report error bars
          (standard deviation across models / runs). Headline
          benchmark numbers in Tables~\ref{tab:binary_results} and
          \ref{tab:multiclass_results} are deterministic at $T=0$.

\item {\bf Experiments compute resources}
    \item[] Question: For each experiment, does the paper provide sufficient information on the computer resources (type of compute workers, memory, time of execution) needed to reproduce the experiments?
    \item[] Answer: \answerYes{}
    \item[] Justification: Appendix~\ref{app:compute}
          itemises GPU-hours per experiment phase (binary,
          multiclass, description, consistency, hidden-state
          extraction) and total ($\sim$845 GPU-hours on RTX 5090s).

\item {\bf Code of ethics}
    \item[] Question: Does the research conducted in the paper conform, in every respect, with the NeurIPS Code of Ethics?
    \item[] Answer: \answerYes{}
    \item[] Justification: We have read the NeurIPS Code
          of Ethics and confirm full conformance. All data was
          collected from publicly licensed sources; faces and license
          plates were screened/removed; no personally identifiable
          information appears in the dataset.

\item {\bf Broader impacts}
    \item[] Question: Does the paper discuss both potential positive societal impacts and negative societal impacts of the work performed?
    \item[] Answer: \answerYes{}
    \item[] Justification: Section~\ref{sec:discussion}
          contains an explicit broader-impacts subsection.
          Positive: the benchmark accelerates safer VLM-based
          perception research. Negative: misinterpretation as an
          endorsement to deploy current VLMs in safety-critical AD;
          we explicitly counsel against this use in the dataset
          documentation.

\item {\bf Safeguards}
    \item[] Question: Does the paper describe safeguards that have been put in place for responsible release of data or models that have a high risk for misuse (e.g., pre-trained language models, image generators, or scraped datasets)?
    \item[] Answer: \answerYes{}
    \item[] Justification: We release no new model weights,
          only image data and annotations. The dataset documentation
          (datasheet) explicitly warns against safety-critical
          deployment based on the released models' performance. All
          potentially-identifiable visual content (faces, plates) was
          screened before release.

\item {\bf Licenses for existing assets}
    \item[] Question: Are the creators or original owners of assets (e.g., code, data, models), used in the paper, properly credited and are the license and terms of use explicitly mentioned and properly respected?
    \item[] Answer: \answerYes{}
    \item[] Justification: All evaluated VLMs and CLIP/SigLIP
          variants are cited with their canonical references. Image
          source datasets and their licenses are listed in
          Appendix~\ref{app:license}; all sources are licensed in a
          manner compatible with CC-BY redistribution.

\item {\bf New assets}
    \item[] Question: Are new assets introduced in the paper well documented and is the documentation provided alongside the assets?
    \item[] Answer: \answerYes{}
    \item[] Justification: \bench is released with a
          datasheet~\citep{datasheets} covering motivation,
          composition, collection process, recommended uses,
          maintenance plan, and known limitations.

\item {\bf Crowdsourcing and research with human subjects}
    \item[] Question: For crowdsourcing experiments and research with human subjects, does the paper include the full text of instructions given to participants and screenshots, if applicable, as well as details about compensation (if any)?
    \item[] Answer: \answerNA{}
    \item[] Justification: No crowdsourced annotation was
          used. Annotations were generated by GPT-4o, validated by
          Claude Sonnet 4.6, and manually adjudicated by the authors.

\item {\bf Institutional review board (IRB) approvals or equivalent for research with human subjects}
    \item[] Question: Does the paper describe potential risks incurred by study participants, whether such risks were disclosed to the subjects, and whether Institutional Review Board (IRB) approvals (or an equivalent approval/review based on the requirements of your country or institution) were obtained?
    \item[] Answer: \answerNA{}
    \item[] Justification: The work does not involve human
          subjects research.

\item {\bf Declaration of LLM usage}
    \item[] Question: Does the paper describe the usage of LLMs if it is an important, original, or non-standard component of the core methods in this research?
    \item[] Answer: \answerYes{}
    \item[] Justification: Appendix~\ref{app:llm_usage}
          declares the use of GPT-4o for initial annotation
          generation and Claude Sonnet 4.6 for cross-validation;
          both are described in detail in
          Sec.~\ref{sec:dataset} and
          Appendix~\ref{app:validation}.

\end{enumerate}


\end{document}
